% Pursuing Stacks (cotes 134-2 to 134-8): archive cote 134-7, site page 48 (PDF page 49), typescript page 490.
% typescript. Draft compared: ps6.tex:2369-2431.
% AI-transcribed by Claude Opus 5.5 (claude-opus-5-5, adaptive thinking, effort high), Message Batches API, batch ps3 (msgbatch_01XdFmuVVBAcgc7EFxdtCqsj), 2026-11-12; not yet checked by eye.
\scanpage{134-7/48}
\mspage{490}

7.9.

116.   For ten days I hav'nt written any notes,  and the time when I
stopped looks a lot more remote still.  For two days still after I last
wrote on the notes,  I kept pondering about schematization of homotopy
types  -  it were rather lively days\over{.}{,}  \over{F}{f}irst day I found the amazingly
simple description of the homotopy groups in the schematic set-up\hand{,} which
got me quite excited;  next day\hand{,} from a phone call to Illusie\hand{,} it turned
out that the key assumption in all my ponderings on schematization,
namely the ``abelianization theorem'' (sic) asserting isomorphism between
discrete and schematic (namely quasi-coherent) \xout{co}homology (or equivalent-%
ly,  cohomology) invariants,  was definitely \uline{false};  consequently,  the
canonical functor from schematic to discrete homotopy types turns out
definitely not to be an equivalence.  This completely overthrows the
idyllic picture in my mind about the relationship between schematic  and
discrete homotopy types  -  but the reflection on schematic homotopy types
``in their own right'' had by then proceeded far enough, the very day befo-%
re,\hand{so} that my faith in the relevance of schematic homotopy types was'nt
seriously shaken  -  rather,  I got excited at drawing a systematic
``bilan'' from the evidence \inserted{now} \deleted{at} hand,about the prospects of developping a
theory of ``schematic homotopy types'' satisfying some basic formal proper-%
ties,  whether or not such theory \over{\unsure{is}}{be} based on semisimplicial unipotent
bundles as \deleted{objects} models,  or \inserted{on} any other kind of models making a sense
over arbitrary ground rings.  Before reverting to a review of the more
formal properties of abeli\over{\unsure{a}}{a}nization in the context of the basic modelizer
(Cat),  I would like still to write down with some care what had thus
come up monday and tuesday last week...

The next days I felt a great fatigue in all my body,  and I then
stopped (till yesterday) any involvment in mathematical reflection.  I
am glad I followed this time the hint that had come to me through my
body,  rather than brush it aside and go on rushing ahead with the work
I was so intensely involved in,  as had been a rule in my life for \xout{\illegible{1}}
many years.  This time,  I understood that the reluctance of the body
to follow that forward rush,  even though I was taking good care of
myself with sleep and food,  had strong reasons,  which had nothing to do
with neither sleep or food or with my general way of life.  Rather,
during the weeks before and also during those very days,  a number of
things had occured in my life,  not all visibly related and of differing
weight and magnitude,  to none of which I had really devoted serious
reflection,  nor even a minimum of time and attention \inserted{needed} for giving me a

\margin{48}
